\documentclass[grau=mestrado, openright]{ppgmat-ime-dissertacao}
\usepackage{dmat-teoremas}
\usepackage{mat-macros}
% Se você não usar subfiles, pode remover o pacote e inserir cada seção do seu trabalho com \include{}. Remova também a linha
\usepackage{subfiles}
\usepackage{lipsum}
\addbibresource{\subfix{bibliografia.bib}} % PARA USO COM O SUBFILES. Caso não for utilizar o subfiles, remova essa linha e descomente a próxima
% \addbibresource{bibliografia.bib}

% ============================
% CONFIGURAÇÕES DO ZREF-CLEVER
% Por padrão utilizamos a opção 'capfirst'. Caso queira alterar o comportamento do comando, confira a documentação do pacote zref-clever para mais informações
% Por exemplo podemos remover o nome da referência. Por exemplo 'Definição 2.3.1' vira '2.3.1'.
% \zcsetup{noname}
% ============================


% Insira as informações de autor, título do trabalho, data, orientador e coorientador (se houver)
\autor{Fulano da Silva}
\titulo{Modelo de dissertação do PPGMAT-IME}
\data{2026}
\orientador{Nome do orientador}
% \coorientador{Nome do coorientador}

% ============================
% OPCIONAL
% ============================
% As informações abaixos já são preenchidas automaticamente pelo modelo. Caso queira alterar alguma delas, basta descomentar o código abaixo.
% \instituicao{%
%   Universidade Federal de Uberlândia -- UFU
%   \par
%   Instituto de Matemática e Estatística
%   \par
%   Graduação}
% \local{Uberlândia - MG}
% O preambulo deve conter o tipo do trabalho, o objetivo, 
% o nome da instituição e a área de concentração 
% \preambulo{Monografia apresentada ao Instituto de Matemática e Estatística, UFU, como requisito para obtenção do título de Bacharel em Matemática,
% sob orientação do Prof. Dr. Orientador da Silva.}
\selectlanguage{brazil}

\begin{document}

% INÍCIO DOS ELEMENTOS PRÉ-TEXTUAIS
\frenchspacing % Espaçamento requerido pela Norma
\pretextual % Marca o início dos elementos pré-textuais a fim de contagem de páginas

\imprimircapa
\imprimirfolhaderosto* % mantenha o asterisco para não alterar a contagem de páginas

% A folha de aprovação deve ser inserida após a defesa. 
% Geralmente você obtém um pdf gerado pelo SEI e o inclui no trabalho pelo comando abaixo.
% \begin{folhadeaprovacao}
% \includepdf{folhadeaprovacao_final.pdf}
% \end{folhadeaprovacao}

\begin{agradecimentos}
    Agradeço...
\end{agradecimentos}

\setlength{\absparsep}{18pt} % ajusta o espaçamento dos parágrafos do resumo
\begin{resumo}
    Este é um modelo baseado no abntex2.

    \vspace{\onelineskip}
    \noindent
    \textbf{Palavras-chave}: Modelo. ABNT. Matemática.
\end{resumo}

\begin{resumo}[Abstract]
    \begin{otherlanguage*}{english}
        This is a template based on abntex2

    \vspace{\onelineskip}
    \noindent
    \textbf{Keywords}: Template. ABNT. Mathematics.

    \end{otherlanguage*}

\end{resumo}
% Sumário
\tableofcontents*
\cleardoublepage


% INÍCIO DOS ELEMENTOS TEXTUAIS
\textual

\subfile{secoes/introducao.tex} % Obrigatório

% --------------
% DESENVOLVIMENTO
% --------------

\subfile{secoes/preliminares.tex}
\subfile{secoes/capitulo1.tex}



% =====================================================
% PÓS TEXTUAL
% =====================================================
\postextual
\printbibliography

\end{document}

